\documentclass[10pt,a4paper]{amsart}
\usepackage[margin=19mm]{geometry}
\usepackage{amsmath,amssymb,mathtools}
\usepackage[scr=rsfs,cal=euler]{mathalfa}
\usepackage{xfrac}
\usepackage[unicode,hidelinks,bookmarks=false]{hyperref}
\newcommand{\ie}{i.e.\ }
\newcommand{\cf}{cf.\ }
\newcommand{\resp}{respectively\ }
\newcommand{\Subsection}{\subsection}
\providecommand{\nobreakdash}{\nobreak\hbox{-}}
\setlength{\emergencystretch}{2em}
\multlinegap0pt
\usepackage{bm}
\usepackage{bbm}
\usepackage{dsfont}
\usepackage{xspace}
\usepackage{cancel}
\usepackage{mathtools}
\usepackage{amsthm}
\makeatletter
\@ifundefined{lemma}{\newtheorem{lemma}{Lemma}}{}
\makeatother
\newcommand{\sn}{\mathsf{n}}
\newcommand{\sr}{\mathsf{r}}
\title[Uniform Approximation of Eigenproblems of a Large-Scale Para]{Uniform Approximation of Eigenproblems of a Large-Scale Parameter-Dependent Hermitian Matrix}
\author{Mattia Manucci, Emre Mengi, Nicola Guglielmi}
\date{Source: arXiv:2409.05791v4 (CC BY 4.0)}
\begin{document}
\maketitle

\noindent\emph{Source and licence.} Uniform Approximation of Eigenproblems of a Large-Scale Parameter-Dependent Hermitian Matrix. Version arXiv:2409.05791v4 (CC BY 4.0), distributed under the Creative Commons Attribution 4.0 International licence, as recorded in the arXiv licence metadata for this exact version and shown on the abstract page https://arxiv.org/abs/2409.05791v4. Full rights notice: licenses/arxiv-2409.05791-CC-BY-4.0.txt.\par\smallskip

\noindent\textit{Editorial scope.} Counted unit: the Lemma whose source label is \texttt{lem:props\_eta} in the article (source0.tex lines 845--876, source label \texttt{lem:props\_eta}), delivered together with the article's own complete original proof of it (source0.tex lines 878--1010, one \texttt{proof} environment of 133 lines). No mathematics is written, repaired or re-derived: statement and proof are the article's own text line for line. Required context retained uncounted: eq:beta\_i (lines 770--781); eq:LP (lines 820--826). Every definition, display and label of those lines is retained unchanged. Omitted from this package and not redistributed: 1--769, 782--819, 827--844, 877--877, 1011--10034. Nothing inside a retained excerpt is omitted or altered: every retained body is delivered exactly as the article's lines are written, so no omission receipt of any kind accompanies this package. Citations: the retained excerpts carry no citation command at all, so no bibliography entry of the article is transcribed and no reference is redistributed. Reference adaptation: none was needed and none was made; every reference inside the delivered unit resolves inside it, so no unresolved reference or citation is produced. Printed slips kept literal: no printed word, symbol, hypothesis or number of the counted statement or of its proof was altered. Layout and class: the article's own class is \texttt{ima-authoring-template} with packages including graphicx, subfigure, color, circuitikz, tikz, comment, pgfplots, caption, xspace, mathtools, cleveref, algorithm, algpseudocode, booktabs; the wrapper is the lane's pinned \texttt{amsart} editorial wrapper at 10pt on A4, which restates the theorem-like environments the retained text needs and adds the article's own macro definitions that the retained text needs (\texttt{\textbackslash sn}, \texttt{\textbackslash sr}), with extra wrapper packages (bm, bbm, dsfont, xspace, cancel, mathtools). The delivered numbering is the unit's own. Assets: no retained span contains a figure environment, a graphics-inclusion command, a PDF-page inclusion, a table or a TikZ picture; no image file is copied into this package, the package declares no asset dependency and no image file is redistributed with it. Licence: version arXiv:2409.05791v4 of this article is distributed under the Creative Commons Attribution 4.0 International licence, as recorded in the arXiv licence metadata for this exact version and shown on the abstract page https://arxiv.org/abs/2409.05791v4. A full rights notice is delivered at licenses/arxiv-2409.05791-CC-BY-4.0.txt. Characters: every retained excerpt is pure ASCII, so the delivered engine, XeLaTeX with its default Latin Modern text font, reports no missing character.
\begin{equation}\label{eq:beta_i}
	\begin{split}
		& \lambda^{{\mathcal U}_j^{\bot}(\mu)}_1(\mu_i)
		\;\geq\;
		\lambda^{(i)}_1 + \beta^{(i,j)}(\mu) \,, \quad \text{where} \\[0.7em]
		& \beta^{(i,j)}(\mu) \;:=\;
		\lambda_{\min} \left( 
			\Lambda^{(i)} - \lambda^{(i)}_1 I_\ell
			- [V^{(i)}]^* U_j(\mu) U_j(\mu)^* V^{(i)} (\Lambda^{(i)} - \lambda^{(i)}_{\ell+1} I_\ell)
		\right) \,,
	\end{split}
\end{equation}

\begin{equation}\label{eq:LP}
	\lambda^{{\mathcal U}_j^{\bot}(\mu)}_1(\mu)
	\;\geq\;
	\eta_\ast^{(j)}(\mu)
	\;:=\;
	\min \left\{ \theta(\mu)^T y \mid y \in \mathcal{Y}^{(j)}_{\mathrm{LB}}(\mu) \right\} .
\end{equation}

\begin{lemma}\label{lem:props_eta}
The following assertions hold for $\eta^{(j)}_\ast(\mu)$ and $\beta^{(i,j)}(\mu)$ defined as in~\eqref{eq:LP} and~\eqref{eq:beta_i}, respectively.
\begin{enumerate}
	\item\label{item1}
	$\beta^{(i,j)}(\mu) \geq 0 \:$ for every $\mu \in {\mathcal D}$ and $i \in \{ 1, \dots , j \}$.
	
	\item\label{item2}
	$\eta^{(j)}_\ast(\mu_i) \geq \lambda^{(i)}_1 \,$ for every $i \in \{ 1, \dots , j \}$.
	
	\item\label{item3}
	If $\lambda^{(i)}_1$ is a simple eigenvalue of $A(\mu^{(i)})$, then
	$\beta^{(i,j)}(\mu_i) > 0 \: $ for every $i \in  \{ 1, \dots , j \}$.
	
	\item\label{item4}
	If $\lambda^{(i)}_1$ is a simple eigenvalue of $A(\mu^{(i)})$, then
	$\eta^{(j)}_\ast(\mu_i) > \lambda^{(i)}_1 \,$ for every $i \in \{ 1, \dots , j \}$.
	
	\item
	For every $i \in \{ 1, \dots , j \}$,
	if $\lambda^{(i)}_{\sr +1} > \lambda^{(i)}_{\sr}$, then 
	\begin{equation}\label{eq:beta_lb}
		\beta^{(i,j)}(\mu_i) = \lambda^{(i)}_{\sr+1} - \lambda^{(i)}_1 > 0.
	\end{equation}

	\item \label{item6}
	For every $i \in \{ 1, \dots , j \}$,
	if $\lambda^{(i)}_{\sr +1} > \lambda^{(i)}_{\sr}$, then
	\begin{equation}\label{eq:lbound_eta}
		\eta^{(j)}_\ast(\mu_i) \;\geq\; \lambda^{(i)}_{\sr +1}.
	\end{equation}
\end{enumerate}
\end{lemma}

\begin{proof}. We proceed point-by-point.
\begin{enumerate}
	\item 
	Observe that $\beta^{(i,j)}(\mu)$ is the smallest eigenvalue of
	\[
		M^{(i)}_j(\mu) \;\vcentcolon=\;
		(\Lambda^{(i)} - \lambda^{(i)}_1 I_\ell) +
		B^{(i)}_j(\mu) B^{(i)}_j(\mu)^\ast (\lambda^{(i)}_{\ell+1} I_\ell - \Lambda^{(i)}),
	\]
	where $B^{(i)}_j(\mu) = [V^{(i)}]^\ast U_j(\mu)$. Moreover, $M^{(i)}_j(\mu)$ is similar to the Hermitian matrix
	\begin{equation}\label{eq:def_tilM}
	\begin{split}
		\widetilde{M}^{(i)}_j(\mu) \;\vcentcolon=\; &
		(\lambda^{(i)}_{\ell+1} I_\ell - \Lambda^{(i)})^{1/2} 
			M^{(i)}_j(\mu)
		(\lambda^{(i)}_{\ell+1} I_\ell - \Lambda^{(i)})^{-1/2} \\[0.5em]
		= \; &
		(\Lambda^{(i)} - \lambda^{(i)}_1 I_\ell) + 
		(\lambda^{(i)}_{\ell+1} I_\ell - \Lambda^{(i)})^{1/2} 
		B^{(i)}_j(\mu) B^{(i)}_j(\mu)^\ast
		(\lambda^{(i)}_{\ell+1} I_\ell - \Lambda^{(i)})^{1/2},
	\end{split}
	\end{equation}
	which is Hermitian positive semidefinite. Hence, $\beta^{(i,j)}(\mu) \geq 0$.

	\item 
	This follows from
	\begin{equation}\label{eq:eta_vs_l1}
		\eta^{(j)}_\ast(\mu_i) = \theta(\mu_i)^T y^{(j)}(\mu_i) 
		\;\geq\; \lambda^{(i)}_1 + \beta^{(i,j)}(\mu_i)
		\;\geq\; \lambda^{(i)}_1,
	\end{equation}
	where the first inequality is due to the definition of $\mathcal{Y}^{(j)}_{\mathrm{LB}}(\mu_i)$
	and $y^{(j)}(\mu_i) \in \mathcal{Y}^{(j)}_{\mathrm{LB}}(\mu_i)$, while the second 
	inequality is due to part~\ref{item1}.

	\item We proceed as in part \ref{item1}. Now the first column of $U_j(\mu_i)$ is an
		eigenvector of $A(\mu_i)$ corresponding to $\lambda^{(i)}_1$. By simplicity
		assumption on $\lambda^{(i)}_1$, this first column is $c v^{(i)}_1$ for some
		$c \in {\mathbb C}$ such that $|c| = 1$, where the eigenvector $v^{(i)}_1$
		is the first column of $V^{(i)}$. By the orthonormality of the columns of $U_j(\mu_i)$
		and $V^{(i)}$, the first column and row of $B^{(i)}_j(\mu_i) = [ V^{(i)} ]^\ast U_j(\mu_i)$
		must be zero except the $(1,1)$ entry which is $c$. The same holds for 
		$B^{(i)}_j(\mu_i) B^{(i)}_j(\mu_i)^\ast$ with the $(1,1)$ entry equal to $| c |^2$, also for 
        \[C \;\vcentcolon=\; (\lambda^{(i)}_{\ell+1} I_\ell - \Lambda^{(i)})^{1/2} 
					B^{(i)}_j(\mu) B^{(i)}_j(\mu)^\ast (\lambda^{(i)}_{\ell+1} I_\ell - \Lambda^{(i)})^{1/2}
                    \]
		with $(1,1)$ entry equal to $|c|^2 (\lambda^{(i)}_{\ell+1} - \lambda^{(i)}_1)$. Now consider $z^\ast \widetilde{M}^{(i)}_j(\mu_i) z$ for any nonzero $z \in {\mathbb C}^{\ell}$,
		where $\widetilde{M}^{(i)}_j(\mu_i)$ is defined as in \eqref{eq:def_tilM}.
		If the first entry of $z$, say $z_1$, is not zero, then letting $\widetilde{z} \in {\mathbb C}^{\ell-1}$ 
		the vector formed of the remaining entries of $z$ excluding its first entry, we have
		\[
			\hskip 5ex
			z^\ast  C z = | z_1 |^2 |c|^2 (\lambda^{(i)}_{\ell+1} - \lambda^{(i)}_1)
					\;	+	\;
							\widetilde{z}^\ast C(2:\ell,2:\ell) \widetilde{z}
					\;\; \geq \;\;	| z_1 |^2 |c|^2 (\lambda^{(i)}_{\ell+1} - \lambda^{(i)}_1) \;\; > \;\; 0,
		\]
		so $z^\ast \widetilde{M}^{(i)}_j(\mu_i) z > 0$. If $z_1=0$, at least one
		of the remaining entries of $z$ is not zero, so
		\[
			z^{\, \ast}
			(\Lambda^{(i)} - \lambda^{(i)}_1 I_\ell)
			z	\;	>	\;	0 \,	,
		\]
		and again $z^\ast \widetilde{M}^{(i)}_j(\mu_i) z > 0$.
		This implies that the smallest eigenvalues of $\widetilde{M}^{(i)}_j(\mu_i)$
		and $M^{(i)}_j(\mu_i)$ are positive, so $\beta^{(i,j)}(\mu_i) > 0$.			
		\smallskip					
		\item 
		This follows from a line of reasoning similar to part \ref{item2}. Specifically,
		(\ref{eq:eta_vs_l1}) holds, but now the last inequality in (\ref{eq:eta_vs_l1}) 
		is satisfied strictly as $\beta^{(i,j)}(\mu_i) > 0$ from part \ref{item3}.		
		\smallskip
		\item Due to the assumption $\: \lambda^{(i)}_{\sr +1} > \lambda^{(i)}_{\sr} \:$, the columns of $U_j(\mu_i)$ 
		and the first $\sr$ columns of $V^{(i)}$ form orthonormal bases for the same invariant 
		subspace of $A(\mu_i)$, namely
		$\text{Null}(A(\mu_i) - \lambda^{(i)}_1 I) \oplus \dots \oplus \text{Null}(A(\mu_i) - \lambda^{(i)}_{\sr} I)$.
		Hence, there is an $\sr \times \sr$ unitary matrix $Q$ such that
		$U_j(\mu_i)	\;	=	\;	 \, V^{(i)}(1:\sn ,  1:\sr) \, Q \:$.
		Now let us first suppose $\ell > \sr$.
		By the orthonormality of the columns of $V^{(i)}$, we have
		\[
			B^{(i)}_j(\mu_i)
				\;	:=	\;
				[ V^{(i)} ]^\ast U_j(\mu_i)
				\;	=	\;
				[ V^{(i)} ]^\ast  V^{(i)}(1:\sn ,  1:\sr) \, Q
				\;	=	\;
			\left[
				\begin{array}{c}
					Q	\\
					0
				\end{array}
			\right]	\:	,
		\]
		which in turn implies
		\begin{align*}
			\begin{aligned}
			C \; 	\vcentcolon= & \; (\lambda^{(i)}_{\ell+1} I_\ell - \Lambda^{(i)})^{1/2} 
				B^{(i)}_j(\mu_i) B^{(i)}_j(\mu_i)^\ast (\lambda^{(i)}_{\ell+1} I_\ell - \Lambda^{(i)})^{1/2}% \\	
			\;	= 	% &\;
	\begin{bmatrix}
					\lambda^{(i)}_{\ell+1} I_{\sr} - \Lambda^{(i)}(1:\sr,1:\sr)	 	&	\;\; 0		\\
						0	&	\;\; 0
				\end{bmatrix}	\:	
			\end{aligned}
		\end{align*}
		so that, recalling (\ref{eq:def_tilM}),
		\begin{equation*}
		\begin{split}
			\widetilde{M}^{(i)}_j(\mu_i)	\;	&	=	\;
			  (\Lambda^{(i)} - \lambda^{(i)}_1 I_\ell)	+	C\;	=\;
				\begin{bmatrix}
					(\lambda^{(i)}_{\ell+1} - \lambda^{(i)}_1)  I_\sr 	 	&	 0		\\
						0	&	 \Lambda^{(i)}(\sr+1:\ell, \sr+1:\ell) - \lambda^{(i)}_1 I_{\ell - \sr}
				\end{bmatrix}	\:	.
		\end{split}
		\end{equation*}	
		Hence, $\beta^{(i,j)}(\mu_i)$, that is the smallest eigenvalue of $M^{(i)}_j(\mu_i)$,
		is also the smallest eigenvalue of $\widetilde{M}^{(i)}_j(\mu)$, which is $\lambda^{(i)}_{\sr +1} - \lambda^{(i)}_1$. If $\ell = \sr$, following the steps of the derivation above, we have $B^{(i)}_j(\mu_i) = Q$,
		$C = \lambda^{(i)}_{\sr+1} I_{\sr} - \Lambda^{(i)}$, 
		$\widetilde{M}^{(i)}_j(\mu) = (\lambda^{(i)}_{\sr+1} - \lambda^{(i)}_1)  I_\sr$, 
		so again $\beta^{(i,j)}(\mu_i) = \lambda^{(i)}_{\sr+1} - \lambda^{(i)}_1$. 
		\smallskip
		\item
		This follows from arguments similar to those used in part \ref{item2}. In particular,
		\eqref{eq:eta_vs_l1} holds, but, using \eqref{eq:beta_lb}, the last inequality in \eqref{eq:eta_vs_l1} can be replaced by $\lambda^{(i)}_1 + \beta^{(i,j)}(\mu_i)  
						\: = \: 
			\lambda^{(i)}_{\sr+1}$.
	\end{enumerate}
        This concludes the proof.
    \end{proof}

\end{document}
